\documentclass[10pt,twocolumn]{article}
\usepackage[margin=0.75in,columnsep=0.28in]{geometry}
\usepackage[T1]{fontenc}
\usepackage{newtxtext,newtxmath}
\usepackage{graphicx,booktabs,amsmath,xcolor,microtype}
\usepackage[hidelinks]{hyperref}
\usepackage[font=small,labelfont=bf]{caption}
\usepackage{enumitem}
\setlist{nosep,leftmargin=1.2em}
\definecolor{gold}{HTML}{B8862E}
\newcommand{\tok}[1]{\texttt{#1}}

\title{\Large\bfseries One Letter: Trace-Grounded Explanatory Film\\of a Single Sampling Step in a 4{,}192-Parameter GPT}
\author{Ziding Zhang\\\small\url{https://zidingz.com/one-letter/}}
\date{\small Working draft, 14 November 2026. Not peer reviewed; not submitted.}

\begin{document}
\maketitle

\begin{abstract}
Explanatory animations of neural networks usually show illustrative numbers chosen by the animator.
We describe a 34-second film that shows one real next-token prediction inside
\emph{microgpt}, a 4{,}192-parameter character-level GPT written in pure Python,
in which every number on screen is read from a recorded execution trace.
Three constraints make the film checkable rather than merely plausible.
(1)~The trace is recorded by an instrumented copy of the model and is accepted only if the run reproduces the
unmodified program's 20 sampled names exactly.
(2)~The depicted step is chosen by a rule fixed before the trace existed, not by taste;
the rule selected the step \tok{ka}\,$\to$\,\tok{m}, where the model sampled a letter of probability 0.027
over a favourite of 0.373.
(3)~Picture and sound are pure functions of one timeline, so their synchrony is verified by measurement:
across 53 audio cues the worst onset error is 3\,ms and every picture event lands within one frame.
We report the pipeline, the numbers it surfaced (in this run, 67 of 118 sampled letters were not the model's top choice),
and its main limitation: we have not yet measured whether viewers learn anything from it, and we describe a study designed to do so.
\end{abstract}

\section{Introduction}
A language model does not output a word; it outputs a probability distribution and then draws from it.
This last step is easy to state and easy to leave out of explanations, which tend to stop at the
``most likely next token.''
Showing it honestly requires a model small enough that \emph{every} intermediate value fits on a screen,
and a production method that prevents the explainer from quietly substituting nicer numbers.

Karpathy's microgpt~\cite{karpathy2026microgpt} provides the first: a complete GPT (tokeniser, scalar autograd,
one transformer layer~\cite{vaswani2017attention}, Adam, sampling) in about 200 lines of dependency-free Python, trained on
${\sim}$32k first names~\cite{makemore}. This paper is about the second.
We contribute:
\begin{itemize}
\item a \textbf{trace contract} with an exact-reproduction acceptance test (\S\ref{sec:trace});
\item a \textbf{pre-committed selection rule} for the depicted example, with a documented revision (\S\ref{sec:select});
\item a \textbf{seekable, deterministic renderer} in which sound and picture share one clock, including a
time-warp for reading pauses that preserves sync by construction (\S\ref{sec:render});
\item \textbf{machine checks} for numerical provenance and audiovisual sync (\S\ref{sec:verify}).
\end{itemize}

\begin{figure*}[t]
\centering
\newcommand{\fr}[2]{\begin{minipage}[t]{0.245\textwidth}\centering\includegraphics[width=\linewidth]{figs/s16-#1.jpg}\\[-1pt]{\scriptsize #2}\end{minipage}}
\fr{a}{(a) the sampled name}\hfill\fr{b}{(b) rewind: 27 options}\hfill\fr{c}{(c) 16 numbers per letter}\hfill\fr{d}{(d) four attention heads}\\[4pt]
\fr{e}{(e) MLP: 7 of 64 active}\hfill\fr{f}{(f) 27 probabilities}\hfill\fr{g}{(g) the die lands on \tok{m}}\hfill\fr{h}{(h) all 20 names; dots mark non-favourites}
\caption{Frames from the 16:9 cut, extracted from the delivered video. Every value shown is read from the recorded trace.}
\label{fig:frames}
\end{figure*}

\section{Related work}
Interactive articles such as Distill's treatment of attention~\cite{olah2016attention} and
tools such as Transformer Explainer~\cite{cho2024transformer} and Bycroft's 3D LLM
visualisation~\cite{bycroft2023llmviz} let readers inspect a running model; they favour exploration over a fixed narrative.
Programmatic animation in the style of Manim~\cite{manim} produces narrative explanations whose
values are usually typed by the author. Segel and Heer~\cite{segel2010narrative} describe the
author-driven end of narrative visualisation, which a short film occupies.
Our work sits between these: a fixed, author-driven film, but one whose numbers are bound to a recorded run
in the way an interactive tool's are. Holtzman et al.~\cite{holtzman2020degeneration} motivate why sampling,
rather than argmax decoding, is the behaviour worth showing.

\section{System}
\subsection{Model and run}
Table~\ref{tab:model} lists the only model facts the film may state.
The run uses the program's own seed and its default inference: 20 samples at temperature $\tau=0.5$,
each ending when the model emits the boundary token \textsc{bos}.

\begin{table}[t]
\centering\small
\begin{tabular}{@{}ll@{}}
\toprule
Training data & ${\sim}$32k names (makemore \tok{names.txt}) \\
Vocabulary & 27 tokens: 26 letters + \textsc{bos} \\
Parameters & 4{,}192 \\
Shape & 1 layer, 4 heads, $d=16$, MLP $16{\to}64{\to}16$ \\
Nonlinearity & ReLU \\
Training & 1{,}000 Adam steps, seed 42, final loss 2.6497 \\
Sampling & $\tau = 0.5$, 20 samples \\
\bottomrule
\end{tabular}
\caption{Model facts. Any other number on screen must come from the trace.}
\label{tab:model}
\end{table}

\subsection{Trace contract}\label{sec:trace}
An instrumented copy of the program records, for every sampled position, the token, its token and position
embeddings, the residual stream before and after the block, per-head attention weights,
the 64 MLP activations before and after ReLU, logits, temperature-scaled probabilities, and the sampled token.
Instrumentation is applied as seven text patches at exact anchors; each anchor must match exactly once,
so a change to the source makes the recorder fail rather than drift.
The patches only \emph{read} values; the patched program is otherwise executed unchanged.

\paragraph{Recording the die.}
Python's \tok{random.choices} draws one uniform $u\in[0,1)$, scales it by the total weight $Z$, and
bisects the cumulative weights $C_j=\sum_{i\le j} w_i$, returning the first $j$ with $C_j > uZ$.
To show where the die actually landed we restate this call with identical arithmetic and record $u$.
Because it consumes the generator identically, the run's outputs are unchanged.

\paragraph{Acceptance.} The trace is accepted only if its 20 names equal those printed by the unmodified program.
They do (20/20), the final training loss matches to four decimals, and the source's SHA-256 is stored in the trace.
A rejected trace is re-recorded, never edited.

\subsection{Selecting the depicted step}\label{sec:select}
To prevent cherry-picking, the brief fixed the selection rule before any trace existed:
\emph{the first sample of length $\ge 4$ containing a position whose sampled token is not the argmax.}
Applied to the trace, this rule (v1) selected position 0 of \tok{kamon}, where the prefix is empty;
all four heads then attend with weight 1.0 to the single position, which makes the attention beat meaningless.
We treated this as a defect in the rule, added the condition \emph{position $\ge 2$} (v2), and recorded the change
before rendering. Rule v2 selects position~2 of the same name: prefix \tok{ka}, sampled \tok{m} with $p=0.0271$,
favourite \tok{r} with $p=0.3727$ (Table~\ref{tab:step}). The recorded draw gives $uZ=0.2264$, inside
\tok{m}'s cumulative segment $[0.2098, 0.2369)$.

\begin{table}[t]
\centering\small
\begin{tabular}{@{}lccc@{}}
\toprule
Head & \textsc{bos} & \tok{k} & \tok{a} \\
\midrule
1 & 0.36 & 0.32 & 0.32 \\
2 & 0.25 & 0.52 & 0.24 \\
3 & 0.01 & \textbf{0.98} & 0.01 \\
4 & 0.38 & 0.61 & 0.01 \\
\bottomrule
\end{tabular}

\smallskip\raggedright
Top 5: \tok{r}~.37, \tok{n}~.27, \tok{l}~.07, \tok{y}~.06, \tok{i}~.04.\\
Sampled: \textcolor{gold}{\tok{m}~.03}. Entropy 1.99 nats. MLP: 7 of 64 units active.
\caption{The selected step, \tok{ka}$\,\to\,$?, read from the trace.}
\label{tab:step}
\end{table}

\subsection{Renderer}\label{sec:render}
Each frame is a pure function $\mathrm{frame} = \mathrm{render}(r, \mathit{fmt})$ of real time $r$ and output format,
drawn with Pillow at $2\times$ supersampling. There is no hidden state: springs and eases are closed-form,
and there is no randomness, so any frame can be rendered alone, out of order, in parallel.
The 16:9 and 9:16 cuts are separate layouts on the same timeline, not crops.

\paragraph{Reading holds.} After a first viewing the film read as too fast. Rather than retiming each shot,
we inserted holds $(h_k, d_k)$ at twelve points where the picture is at rest. Shots are authored in
\emph{film time} $t$; real time is
\begin{equation}
r(t) = t + \textstyle\sum_k d_k\,\mathbb{1}[h_k < t],
\end{equation}
and the renderer evaluates the scene at $f(r)$, the inverse of $r$, which is constant during a hold.
Every audio cue is authored in film time and placed at $r(t)$, so the 9.7\,s of holds (24.0\,s $\to$ 33.7\,s)
cannot break synchrony. We checked that every hold starts on a frame with zero mean pixel change.

\paragraph{Sound from the same clock.} All sound is synthesised (sums of decaying sines): a tick per letter,
section downbeats, one tone per attention head scaled by that head's largest weight, a low drone whose
level tracks the entropy of the step being decided, and a struck note when the die lands.
The die marker moves as $x(s) = uZ\cdot E(s)$ with ease-out $E(s)=1-(1-s)^3$ over $s\in[0,1]$.
It clicks as it crosses each boundary $C_j < uZ$; the click times are computed by inverting the ease,
\begin{equation}
s_j = 1 - \left(1 - C_j/(uZ)\right)^{1/3},
\end{equation}
so the twelve clicks before \tok{m} are exact, not placed by ear.


\section{Verification}\label{sec:verify}
\paragraph{Provenance.} The renderer reads values only from the trace and a derived summary of the selected step;
no number is typed into the scene code.
\paragraph{Synchrony.} A script compares the audio and the rendered frames against the 53 cues
without listening. Audio onsets are found as peaks of positive energy flux on a high-passed signal at 1\,ms hops;
slow-attack cues are checked as a level rise instead. For each discrete event it checks that the first frame
at or after the cue differs from the previous frame. It also cross-correlates the muxed file's audio with the source WAV.
\paragraph{Results.} Isolated onsets fall within 3\,ms of their cue (tolerance 10\,ms); every picture event
changes the frame within one frame (${\le}33$\,ms at 30\,fps); the muxed audio offset is 0.0\,ms in both cuts;
0 of 53 checks fail. The check also caught a real defect: a swell under the growing bars was masked by a
downbeat's tail ($-2.6$\,dB) and was cut.

\section{What the trace showed}
The film's closing claim was not planned. Counting across the whole run, 67 of the 118 sampled letters were
not the model's most likely letter. At $\tau = 0.5$ this model leaves the favourite more often than it takes it,
so the featured step is typical rather than a rare accident. We state this only for this run;
other seeds would give other counts.

\section{Planned evaluation}
We have designed, but not yet run, a between-subjects study ($3$ arms, about $110$ participants each) with a pre-test,
an immediate post-test and an optional one-week follow-up.
\textbf{A} is the film as published. \textbf{B} is a favourite-only film, already rendered: the same step \tok{ka}$\,\to\,$? with identical numbers,
taken from \tok{karai}, another name in the same run that reached the same prefix and wrote the favourite \tok{r} ($p=0.37$).
It has the same length, timeline and sound, but no die, no not-the-favourite marks and no closing count,
so it keeps the no-invented-numbers rule while removing only the sampling story.
\textbf{C} is a 120-word text with the same content as A.
The primary outcome is a 9-item sampling-understanding score (e.g.\ \emph{``out of 100 runs, how often would it write a 3\% letter?''},
plus transfer items on temperature and chatbot variability). H1: A\,$>$\,B; H2: A\,$>$\,C (ANCOVA on post-test with pre-test as covariate,
Holm-corrected; $n=100$ per arm gives power $0.80$ for $d=0.4$). H3, exploratory: the A$-$B gap is larger for viewers with little ML experience.
Exclusions, items and analysis will be preregistered before data collection.

\section{Limitations}
\begin{itemize}
\item \textbf{No audience evaluation yet.} Until the study above runs, we cannot claim the film teaches anything.
\item \textbf{One run, one step.} All numbers come from one seed; the film generalises only through the run-wide count.
\item \textbf{Simplified display.} Embedding and activation values are shown as colour on a shared clamp
(98th percentile, square-root curve), so small values are visibly compressed.
\item \textbf{Process.} The pipeline was built by an AI coding agent under staged human approval;
the gates are documented, but taste judgements (sound, typeface) were made by one person.
\end{itemize}

\section{Conclusion}
A tiny model makes it possible to show a language model's choice with no invented numbers.
The cost of that honesty is mostly engineering discipline: an acceptance test for the trace, a selection rule
written before the data, and one clock for picture and sound. The resulting checks are cheap enough to rerun on every render.

\paragraph{Availability.}\raggedright Film: \url{https://zidingz.com/one-letter/}. Source and study films: \url{https://github.com/sartoshi-foot-dao/one-letter}.
Model by Andrej Karpathy~\cite{karpathy2026microgpt}; data from makemore (MIT)~\cite{makemore}.

\paragraph{Acknowledgements.} The renderer, trace recorder and verification scripts were written with Claude (Anthropic).

{\footnotesize
\begin{thebibliography}{9}\setlength{\itemsep}{1pt}
\bibitem{karpathy2026microgpt} A.~Karpathy. microgpt. Blog post and code, Feb.\ 2026. \url{https://karpathy.github.io/2026/02/12/microgpt/}
\bibitem{makemore} A.~Karpathy. makemore. GitHub repository, MIT licence. \url{https://github.com/karpathy/makemore}
\bibitem{olah2016attention} C.~Olah and S.~Carter. Attention and augmented recurrent neural networks. \emph{Distill}, 2016.
\bibitem{cho2024transformer} A.~Cho et al. Transformer Explainer: Interactive learning of text-generative models. arXiv:2408.04619, 2024.
\bibitem{bycroft2023llmviz} B.~Bycroft. LLM Visualization. \url{https://bbycroft.net/llm}, 2023.
\bibitem{manim} G.~Sanderson and the Manim Community. Manim: a mathematical animation engine. Software.
\bibitem{segel2010narrative} E.~Segel and J.~Heer. Narrative visualization: Telling stories with data. \emph{IEEE TVCG} 16(6), 2010.
\bibitem{holtzman2020degeneration} A.~Holtzman, J.~Buys, L.~Du, M.~Forbes, Y.~Choi. The curious case of neural text degeneration. \emph{ICLR}, 2020.
\bibitem{vaswani2017attention} A.~Vaswani et al. Attention is all you need. \emph{NeurIPS}, 2017.
\end{thebibliography}}
\end{document}
